% GENERATED FILE. Edit canonical lessons, then run npm run generate:books.
% canonical-source-hash: fnv1a64:98fa666f652d193d
% canonical-lessons: JA-C80-kousaten, JA-C80-shingou, JA-C80-hougaku, JA-C80-chikamichi, JA-C80-mawari

\chapter{Crossroads, Traffic light, Direction, Shortcut, Surroundings}
\label{ch:ja-crossroads-traffic-light-direction-shortcut-surroundings}
\hlchaptermodality{80}

\begin{chapteropening}
\textbf{By the end of this chapter:} \emph{I can say a crossroads, a traffic light, a direction, a shortcut and the surroundings.}

Complete the last lesson of chapter 80: i can say a crossroads, a traffic light, a direction, a shortcut and the surroundings.
\end{chapteropening}

\section[\texorpdfstring{\ja{こうさてん} (kousaten)}{こうさてん (kousaten)}]{\ja{こうさてん} (kousaten) — a crossroads}

\label{lesson:JA-C80-kousaten}



\begin{quote}
Before the new one: say the Japanese for stairs, then the Japanese for a corridor.
\end{quote}

\subsection*{You'll want to know: \ja{こうさてん}}
\textbf{\ja{こうさてん}} — \emph{kousaten} — \textquotedblleft{}a crossroads\textquotedblright{}.

\begin{grammarlens}[title={What you've built}]
One more word for this chapter.
\end{grammarlens}

\subsection*{Your turn}
\begin{itemize}
\raggedright
  \item \emph{Say it:} \emph{kousaten}
  \item \emph{Say it:} \emph{kousaten}, once more
  \item \emph{Say it:} say \emph{kousaten}
  \item \emph{Recall:} say the Japanese for stairs, then the Japanese for a corridor, then say \emph{kousaten} again
\end{itemize}

\begin{tcolorbox}[breakable,colback=teal!4,colframe=teal!35!black,title={Before you move on}]
What does \ja{こうさてん} mean? (\textquotedblleft{}A crossroads\textquotedblright{}.) Say it once more.
\end{tcolorbox}
\section[\texorpdfstring{\ja{しんごう} (shingou)}{しんごう (shingou)}]{\ja{しんごう} (shingou) — a traffic light}

\label{lesson:JA-C80-shingou}



\begin{quote}
Before the new one: say the Japanese for a corridor, then the Japanese for a crossroads.
\end{quote}

\subsection*{You'll want to know: \ja{しんごう}}
\textbf{\ja{しんごう}} — \emph{shingou} — \textquotedblleft{}a traffic light\textquotedblright{}.

\begin{grammarlens}[title={What you've built}]
One more word for this chapter.
\end{grammarlens}

\subsection*{Your turn}
\begin{itemize}
\raggedright
  \item \emph{Say it:} \emph{shingou}
  \item \emph{Say it:} \emph{shingou}, once more
  \item \emph{Say it:} say \emph{shingou}
  \item \emph{Recall:} say the Japanese for a corridor, then the Japanese for a crossroads, then say \emph{shingou} again
\end{itemize}

\begin{tcolorbox}[breakable,colback=teal!4,colframe=teal!35!black,title={Before you move on}]
What does \ja{しんごう} mean? (\textquotedblleft{}A traffic light\textquotedblright{}.) Say it once more.
\end{tcolorbox}
\section[\texorpdfstring{\ja{ほうがく} (hougaku)}{ほうがく (hougaku)}]{\ja{ほうがく} (hougaku) — a direction}

\label{lesson:JA-C80-hougaku}



\begin{quote}
Before the new one: say the Japanese for a crossroads, then the Japanese for a traffic light.
\end{quote}

\subsection*{You'll want to know: \ja{ほうがく}}
\textbf{\ja{ほうがく}} — \emph{hougaku} — \textquotedblleft{}a direction\textquotedblright{}.

\begin{grammarlens}[title={What you've built}]
One more word for this chapter.
\end{grammarlens}

\subsection*{Your turn}
\begin{itemize}
\raggedright
  \item \emph{Say it:} \emph{hougaku}
  \item \emph{Say it:} \emph{hougaku}, once more
  \item \emph{Say it:} say \emph{hougaku}
  \item \emph{Recall:} say the Japanese for a crossroads, then the Japanese for a traffic light, then say \emph{hougaku} again
\end{itemize}

\begin{tcolorbox}[breakable,colback=teal!4,colframe=teal!35!black,title={Before you move on}]
What does \ja{ほうがく} mean? (\textquotedblleft{}A direction\textquotedblright{}.) Say it once more.
\end{tcolorbox}
\section[\texorpdfstring{\ja{ちかみち} (chikamichi)}{ちかみち (chikamichi)}]{\ja{ちかみち} (chikamichi) — a shortcut}

\label{lesson:JA-C80-chikamichi}



\begin{quote}
Before the new one: say the Japanese for a traffic light, then the Japanese for a direction.
\end{quote}

\subsection*{You'll want to know: \ja{ちかみち}}
\textbf{\ja{ちかみち}} — \emph{chikamichi} — \textquotedblleft{}a shortcut\textquotedblright{}.

\begin{grammarlens}[title={What you've built}]
One more word for this chapter.
\end{grammarlens}

\subsection*{Your turn}
\begin{itemize}
\raggedright
  \item \emph{Say it:} \emph{chikamichi}
  \item \emph{Say it:} \emph{chikamichi}, once more
  \item \emph{Say it:} say \emph{chikamichi}
  \item \emph{Recall:} say the Japanese for a traffic light, then the Japanese for a direction, then say \emph{chikamichi} again
\end{itemize}

\begin{tcolorbox}[breakable,colback=teal!4,colframe=teal!35!black,title={Before you move on}]
What does \ja{ちかみち} mean? (\textquotedblleft{}A shortcut\textquotedblright{}.) Say it once more.
\end{tcolorbox}
\section[\texorpdfstring{\ja{まわり} (mawari)}{まわり (mawari)}]{\ja{まわり} (mawari) — the surroundings}

\label{lesson:JA-C80-mawari}



\begin{quote}
Before the new one: say the Japanese for a direction, then the Japanese for a shortcut.
\end{quote}

\subsection*{You'll want to know: \ja{まわり}}
\textbf{\ja{まわり}} — \emph{mawari} — \textquotedblleft{}the surroundings\textquotedblright{}.

\begin{grammarlens}[title={What you've built}]
One more word for this chapter.
\end{grammarlens}

\subsection*{Your turn}
\begin{itemize}
\raggedright
  \item \emph{Say it:} \emph{mawari}
  \item \emph{Say it:} \emph{mawari}, once more
  \item \emph{Say it:} say \emph{mawari}
  \item \emph{Recall:} say the Japanese for a direction, then the Japanese for a shortcut, then say \emph{mawari} again
\end{itemize}

\begin{tcolorbox}[breakable,colback=teal!4,colframe=teal!35!black,title={Before you move on}]
What does \ja{まわり} mean? (\textquotedblleft{}The surroundings\textquotedblright{}.) Say it once more.
\end{tcolorbox}
